\chapter{Analyse et conception}

\newpage

\section*{Introduction}
\addcontentsline{toc}{section}{Introduction}


\section{Concepts de base}

\subsection{Domain Driven Design}


% Les concepts clés de DDD incluent les entités, les objets de valeur, les événements de domaine, les agrégats, les services, les référentiels et les fabriques.


\subsection{Architecture Microservices}

\subsubsection{Définition}


\subsubsection{Patterns}


\section{Analyse}

\subsection{Analyse de l'existant}


\begin{table}[H]
    \hspace{-1cm}
    \begin{tabular}{ |m{4cm}| m{3cm}| m{10cm} |}
    \hline
     Acteur & Type & Description\\ 
    \hline
    ECAPI Client 
        & Principale
         &
         L'application qui interagit avec l’ECAPI Cart pour l’ajout, la suppression, ou bien la modification des produits. Cela peut être un autre ECAPI, l’application web NC2, ou bien l'application mobile. \\
    \hline
    ECAPI Customer 
    
        & Secondaire 
        &
        L’application responsable de la gestion des clients du système. ECAPI Cart a besoin de cette application pour faire l’authentification de client.
        \\
    \hline
    \end{tabular}
    \caption{Acteurs et leurs rôles dans le système}
\end{table}

\section*{Conclusion}
\addcontentsline{toc}{section}{Conclusion}





